\chapter{Side-Channel Security}
\label{ch:sidechannel}
\chaptermark{Side-Channel Security}

\section{Learning objectives}

Every earlier chapter of this book argued about security \emph{on paper}: a lattice
problem is hard, a forgery reduces to Module-SIS, a distinguisher would break LWE.
That kind of security is necessary. It is also not sufficient. A scheme does not run
on paper; it runs on a physical device that takes time, draws current, radiates,
and can be poked with a laser. If any of those measurable, physical quantities
depends on the secret key, an attacker can read the key off the machine without ever
touching the mathematics. This chapter's warning is blunt:
\emph{the math can be perfect and the deployment still broken.} It opens the final part
of the book, on practice: correct, side-channel-safe implementation is one requirement
among several for getting these schemes into real systems, alongside the protocol
integration and migration that the chapters after it take up.

After finishing this chapter, you should be able to:

\begin{enumerate}
  \item explain why provable, mathematical security is necessary but not sufficient,
        and say precisely what a \emph{side channel} is;
  \item name the main families of implementation attacks -- timing, cache, power,
        electromagnetic, and fault -- and state what each one actually observes;
  \item state the rules of \emph{constant-time} programming and write a constant-time
        conditional select and a constant-time comparison;
  \item connect these ideas to concrete PQC design choices already met in this book:
        CBD instead of Gaussian sampling (\ref{ch:cbd}), Falcon's floating-point
        sampler (\ref{ch:fndsa}), ML-DSA's deterministic abort loop (\ref{ch:mldsa}),
        and the implicit rejection built into ML-KEM decapsulation (\ref{ch:fips203});
  \item describe the higher-order countermeasures -- masking and blinding -- and
        weigh their cost against constant-time coding.
\end{enumerate}

\paragraph{Why this chapter matters.}
Shor's algorithm, the threat that motivated the whole book, attacks the
\emph{structure} of the problem: it factors, it computes discrete logs. The schemes
we built defeat it because no efficient quantum algorithm is known for the lattice
and code problems underneath. A side-channel attack does something completely
different. It ignores the problem structure and instead measures the execution:
how long a decryption took, which cache lines it touched, how much power the chip
drew on each clock cycle. These attacks are often far \emph{cheaper} than the
mathematical ones -- minutes with an oscilloscope rather than aeons on a quantum
computer -- which is exactly why they matter. The concern is not hypothetical for
PQC: several of the schemes in this book have delicate implementation requirements,
and those requirements shaped the standards. It is why ML-KEM samples noise with CBD
rather than a Gaussian (\ref{ch:cbd}), and why FN-DSA (based on Falcon; FIPS~206
remains in development) is especially delicate to deploy safely (\ref{ch:fndsa}).

\section{The attack families}

A side channel is any observable output of the computation other than the intended
result. Below are the families a PQC implementer must keep in mind. Each is sketched
in one paragraph; the common thread is that all of them turn a physical measurement
into information about the secret.

\paragraph{Timing attacks.}
The total time a computation takes depends on the data it processed~\cite{kocher1996}. If that data
includes the secret, the running time leaks it. The textbook example is a
comparison that returns as soon as it finds the first differing byte: an attacker
who can vary an input and time the response learns the secret one byte at a time.
Timing leaks are the most dangerous class because they need no special equipment
and can often be measured \emph{remotely}, over a network.

\paragraph{Cache attacks.}
Modern CPUs cache recently used memory. If a program indexes a table at a
secret-dependent position -- as many Gaussian samplers and cipher S-boxes do -- then
which table entries land in the cache depends on the secret. An attacker sharing the
same machine (a co-resident cloud tenant, say) does not see the values, but can time
its \emph{own} memory accesses to learn which cache lines were touched, using
techniques such as Flush+Reload or Prime+Probe. Secret-dependent memory addresses are
therefore just as leaky as secret-dependent timing.

\paragraph{Power analysis.}
The instantaneous power a chip consumes correlates with the operations it performs
and the bits it moves. \emph{Simple power analysis} (SPA) reads a single trace to
recognize the sequence of operations -- for instance, telling a multiply from a
square in a modular exponentiation, which spells out exponent bits. \emph{Differential
power analysis} (DPA)~\cite{kjj1999} is stronger: it collects many traces and correlates them,
sample by sample, against a hypothesis about a few key bits, pulling the secret out
even when it sits far below the measurement noise.

\paragraph{Electromagnetic analysis.}
A switching transistor radiates. The same information that power analysis reads from
the supply line can be read from the electromagnetic emanations near the chip, often
with a small probe and no electrical contact. EM analysis can also be spatially
localized to a single region of the die, which sometimes defeats countermeasures that
only balance the \emph{global} power draw.

\paragraph{Fault injection.}
The attacks above are passive: they observe. Fault injection is active. By glitching
the supply voltage or clock, or firing a laser or EM pulse at the die, an attacker
induces a controlled error -- a skipped instruction, a flipped bit, a corrupted
value -- and exploits the faulty output. Signature schemes are especially fragile
here: a single faulty signature can, through differential fault analysis, reveal the
signing key, and skipping a bounds check can leak the mask in a scheme like ML-DSA.

\section{Constant-time programming}

The central software defense against secret-dependent \emph{timing and cache} leakage is \textbf{constant-time programming}. It is necessary but not sufficient for physical power or EM attacks. The name is slightly misleading: the
goal is not that the code always takes the same wall-clock time, but that its timing,
memory-access pattern, and operation sequence \emph{do not depend on secret data}.
Variation that depends only on public data (the length of a message, the public
iteration count of a loop) is fine.

The discipline reduces to three rules:

\begin{itemize}
  \item \textbf{No secret-dependent branches.} An \texttt{if} whose condition
        involves a secret takes one path or another, and the two paths almost never
        cost the same. Replace the branch with straight-line arithmetic.
  \item \textbf{No secret-dependent memory indices.} Never use a secret as an array
        index or a jump target. A table lookup at a secret address leaks through the
        cache even if the code has no branches at all.
  \item \textbf{Avoid variable-time instructions on secret data.} Division and modulo,
        and on some processors even multiplication, take a data-dependent number of
        cycles. Early-exit library routines such as \texttt{memcmp} are the same
        hazard in disguise. Use fixed-cost operations, or reduce with Montgomery/Barrett
        arithmetic instead of a hardware divide.
\end{itemize}

Figure~\ref{fig:ct-branch} contrasts the hazard with its fix. On the left, a secret
bit chooses between a short path and a long path; the difference in duration is the
leak. On the right, the same choice is computed with a mask so that every input runs
the identical sequence of operations.

\begin{figure}[H]
\centering
\begin{tikzpicture}[every node/.style={font=\scriptsize}]
  % ---- Left: leaky secret-dependent branch ----
  \node[font=\small\bfseries] at (3.1,4.9) {Secret-dependent branch (leaks)};
  \node[flownode] (inA) at (0,3.3) {input};
  \node[diamond, draw=pqcorange, thick, fill=white, align=center,
        inner sep=1pt] (d) at (1.9,3.3) {secret\\bit?};
  \node[flownode] (t1) at (4.1,4.2) {op};
  \node[flownode] (b1) at (3.3,2.3) {op};
  \node[flownode] (b2) at (4.5,2.3) {op};
  \node[flownode] (b3) at (5.7,2.3) {op};
  \node[keynode]  (outA) at (7.2,3.3) {result};
  \node[font=\tiny, pqcteal] at (2.7,4.15) {bit $=1$};
  \node[font=\tiny, pqcorange] at (2.55,2.62) {bit $=0$};
  \draw[flowarrow] (inA) -- (d);
  \draw[flowarrow] (d) |- (t1);
  \draw[flowarrow] (t1) -| (outA);
  \draw[flowarrow] (d) |- (b1);
  \draw[flowarrow] (b1) -- (b2);
  \draw[flowarrow] (b2) -- (b3);
  \draw[flowarrow] (b3) -| (outA);
  \node[font=\tiny, pqcorange, align=center] at (3.5,1.4)
       {one path is short, the other long:\\ duration reveals the secret bit};

  % ---- Divider ----
  \draw[pqcgray!40, dashed] (-0.9,0.7) -- (8.3,0.7);

  % ---- Bottom: constant-time straight line ----
  \node[font=\small\bfseries] at (3.1,0.1) {Constant-time (timing/cache-safe)};
  \node[flownode] (c0) at (0,-1.1) {input};
  \node[flownode] (c1) at (1.9,-1.1) {op};
  \node[flownode] (c2) at (3.4,-1.1) {op};
  \node[flownode] (c3) at (4.9,-1.1) {op};
  \node[flownode] (c4) at (6.2,-1.1) {mask};
  \node[keynode]  (c5) at (7.7,-1.1) {result};
  \draw[flowarrow] (c0) -- (c1);
  \draw[flowarrow] (c1) -- (c2);
  \draw[flowarrow] (c2) -- (c3);
  \draw[flowarrow] (c3) -- (c4);
  \draw[flowarrow] (c4) -- (c5);
  \node[font=\tiny, pqcteal, align=center] at (3.7,-2.0)
       {every input runs the same operations:\\ duration is independent of the secret};
\end{tikzpicture}
\caption{Top: a branch on a secret bit takes a short or a long path, so the running time and cache trace can carry the secret. Bottom: the same conditional is computed as one fixed sequence of operations ending in a masked select. This removes secret-dependent control flow and memory behaviour, but intermediate values can still correlate with physical power or EM leakage.}
\label{fig:ct-branch}
\end{figure}

The workhorse that removes a secret branch is the \emph{constant-time select}. It
turns the selector bit into an all-ones or all-zero \emph{mask} and uses that mask to
blend the two candidate values with bitwise arithmetic, so both candidates are always
computed and no branch is taken.

\begin{algorithm}[H]
\caption{\textsc{CTSelect} -- constant-time conditional select}
\label{alg:ctselect}
\begin{algorithmic}[1]
\Require Selector bit $s \in \{0,1\}$; machine words $a$ and $b$
\Ensure $a$ if $s=1$, else $b$ -- with no data-dependent branch
\State $mask \gets -\,s$
       \Comment{two's complement: all-ones word if $s=1$, all-zero if $s=0$}
\State \Return $(a \mathbin{\&} mask) \mathbin{|} (b \mathbin{\&} \overline{mask})$
       \Comment{$\overline{mask}$ is the bitwise complement of $mask$}
\end{algorithmic}
\end{algorithm}

The trick is the mask $-s$: negating the bit $1$ in two's complement gives the
all-ones word, and negating $0$ gives the all-zero word. When $s=1$ the first term
passes $a$ and the second term is zeroed; when $s=0$ the roles swap. There is no
branch and no secret-dependent address. The same masking idea gives a constant-time
\emph{equality} test, \textsc{ct\_eq}, which must never return early: it accumulates
the differences across \emph{all} bytes and only then collapses the result to a
single bit.

Listing~\ref{lst:ct-c} shows both patterns in C: the leaky branch, its masked
replacement, and a no-early-exit comparison.

\begin{lstlisting}[language=C,label=lst:ct-c,caption={Leaky versus constant-time patterns in C.}]
#include <stdint.h>
#include <stddef.h>

/* LEAKY: the branch condition is a secret; the two paths differ */
uint32_t leaky_select(int secret_bit, uint32_t a, uint32_t b) {
    if (secret_bit)      /* timing/cache/power all reveal which way */
        return a;
    else
        return b;
}

/* CONSTANT-TIME: no branch, no secret-dependent memory access */
uint32_t ct_select(uint32_t secret_bit, uint32_t a, uint32_t b) {
    uint32_t mask = -secret_bit;   /* 0xFFFFFFFF if 1, 0x00000000 if 0 */
    return (a & mask) | (b & ~mask);
}

/* CONSTANT-TIME equality: 1 if equal, else 0, never exits early */
uint32_t ct_eq(const uint8_t *x, const uint8_t *y, size_t n) {
    uint8_t diff = 0;
    for (size_t i = 0; i < n; i++)
        diff |= x[i] ^ y[i];       /* fold every byte, no break */
    /* top bit of (diff | -diff) is 1 iff diff != 0 */
    uint32_t nz = ((uint32_t)diff | (0u - (uint32_t)diff)) >> 31;
    return nz ^ 1u;                /* invert: equal -> 1 */
}
\end{lstlisting}

\section{PQC-specific concerns}

Constant-time discipline is general, but post-quantum schemes concentrate the risk in
a few specific places. Each of the following ties back to a design choice made in an
earlier chapter.

\paragraph{Sampling: the reason CBD exists.}
The noise in a lattice scheme must be sampled every time a key or ciphertext is made.
A discrete \emph{Gaussian} sampler is the natural mathematical choice, but it is a
side-channel minefield: it typically needs rejection loops (a secret-dependent number
of iterations), table lookups (secret-dependent memory addresses), and floating-point
arithmetic (variable-time on most hardware). ML-KEM samples its small noise
polynomials from the \emph{centered binomial distribution} of
Chapter~\ref{ch:cbd}. ML-DSA uses its own bounded-uniform and rejection-sampling
procedures for secrets and masks; it is not an instance of the ML-KEM CBD sampler.
Both standards avoid a general-purpose discrete-Gaussian sampler, but their sampling
interfaces and side-channel reviews must be considered separately. FN-DSA
(\ref{ch:fndsa}) is the cautionary opposite. Its trapdoor sampler genuinely needs a high-precision Gaussian
over the reals, computed with floating-point FFT, and making \emph{that} constant-time
is the hardest implementation problem among all the standards -- if the sampler's
timing or rounding leaks, the secret basis can be reconstructed. That difficulty,
forward-referenced at the end of Chapter~\ref{ch:fndsa}, is this chapter: it is the
single strongest reason Falcon is the most demanding of the three signature standards
to deploy safely.

\paragraph{Rejection sampling in ML-DSA.}
ML-DSA signing (\ref{ch:mldsa}) also loops: it draws a mask, forms
$\mathbf{z}=\mathbf{y}+c\,\mathbf{s}_1$, and applies rejection conditions involving
secret-dependent quantities, including checks related to $\mathbf{s}_2$ and
$\mathbf{t}_0$. The retry count and signing time therefore must not be declared
harmless merely from distributional intuition. Implementations must follow FIPS~204's
sampling and side-channel design, protect intermediate values, and evaluate the
resulting machine code on the target platform. FIPS~204 supports hedged signing, with
randomness derived from secret $K$, the message representative $\mu$, and input $rnd$;
deterministic operation fixes $rnd$, but does not remove $K$ or the secret-related
rejection conditions.

\paragraph{Decapsulation and the Fujisaki--Okamoto transform.}
The subtlest PQC side channel lives in KEM decapsulation. A raw lattice decryption can
\emph{fail}, and whether it failed depends on the secret key and the (possibly
attacker-chosen) ciphertext. If decapsulation ever behaved differently on success
versus failure -- returned an error, took longer, or drew different power -- the
attacker would gain a \emph{decryption-failure} or \emph{plaintext-checking} oracle,
querying malformed ciphertexts and using the yes/no answers to peel the secret apart.
The Fujisaki--Okamoto (FO) transform~\cite{fo1999}, developed in
Section~\ref{sec:fo}, closes this door with \textbf{implicit
rejection}. Decapsulation re-encrypts the recovered message and compares the result to
the received ciphertext in constant time. On a match it returns the real shared key;
on a mismatch it returns a \emph{pseudorandom} key derived from a secret seed and the
ciphertext -- not a distinguishable ciphertext-validity error. This applies after
basic input and decapsulation-key checks; an API may still reject malformed types,
lengths, or unusable key material. Both cryptographic outcomes produce a key of identical form, and the
final choice is a constant-time select, so the caller cannot tell success from
failure. Figure~\ref{fig:fo-implicit} shows the flow.

\begin{figure}[t]
\centering
\begin{tikzpicture}[every node/.style={font=\scriptsize}]
  \node[flownode] (ct)   at (2,6.1) {ciphertext $ct$};
  \node[flownode] (dec)  at (2,4.9) {$m' = \textsc{Decrypt}(sk, ct)$};
  \node[flownode, align=center] (re) at (2,3.5)
       {$(K',r) = G(m'\,\|\,h)$\\ $ct' = \textsc{Encrypt}(pk, m'; r)$};
  \node[diamond, draw=pqcorange, thick, fill=white, align=center, inner sep=1pt]
       (cmp) at (2,1.7) {$ct' = ct$\\ ?};
  \node[flownode] (kreal) at (6.2,3.5) {real key $K'$};
  \node[flownode, align=center] (kfake) at (6.2,0.4)
       {$\bar K = J(z, ct)$\\ (pseudorandom)};
  \node[flownode, align=center] (sel) at (6.2,1.9)
       {\textsc{CTSelect}\\ $(\text{match},K',\bar K)$};
  \node[keynode] (out) at (9.4,1.9) {shared key $K$};
  \draw[flowarrow] (ct)  -- (dec);
  \draw[flowarrow] (dec) -- (re);
  \draw[flowarrow] (re)  -- (cmp);
  \draw[flowarrow] (re.east) -- (kreal.west);
  \draw[flowarrow] (kreal) -- (sel);
  \draw[flowarrow] (kfake) -- (sel);
  \draw[flowarrow] (cmp) -- node[above,font=\tiny,pqcgray]{match bit} (sel);
  \draw[flowarrow] (sel) -- (out);
  \node[font=\tiny, pqcgray, align=center] (zseed) at (3.7,0.4)
       {$z$: secret\\ reject seed};
  \draw[flowarrow, dashed] (zseed) -- (kfake.west);
\end{tikzpicture}
\caption{FO decapsulation with implicit rejection. The recovered message is
re-encrypted and compared to the received ciphertext in constant time. A match yields
the real key $K'$; a mismatch yields a pseudorandom key $\bar K = J(z,ct)$ derived
from a secret seed $z$. Because the two outputs have identical form and the choice is
a constant-time select, success and failure are indistinguishable to the caller -- no
decryption-failure oracle exists.}
\label{fig:fo-implicit}
\end{figure}

\paragraph{The NTT and polynomial arithmetic.}
The number-theoretic transform and the modular polynomial arithmetic that surround it
run on secret data on every operation. They must be constant-time too: the modular
reductions after each butterfly should use fixed-cost Montgomery or Barrett reduction
rather than a hardware division, and there must be no secret-dependent branch anywhere
in the transform. This is usually easy to achieve -- the NTT is naturally
straight-line -- but it is easy to spoil with a careless conditional subtraction that
runs only ``when the coefficient is too large.''

\section{Countermeasures beyond constant-time}

Constant-time coding primarily addresses timing, cache, and secret-dependent control-flow leakage. Differential power and EM analysis are stronger, because they average many
traces and can pull out a secret that influences the power draw even when the control
flow is perfectly uniform. Defeating those needs a further layer.

\paragraph{Masking.}
The main tool is \textbf{masking}: split every secret value into several random
\emph{shares} whose combination (XOR, for Boolean masking; addition mod $q$, for
arithmetic masking) equals the secret, and compute on the shares so that no single
intermediate value in the whole computation ever equals, or correlates with, the
secret. A first-order masking uses two shares and forces the attacker to combine
measurements from two points at once; an order-$d$ masking uses $d+1$ shares. The
security grows with $d$, but so does the cost -- the number of shares grows linearly
while the number of secure operations and the fresh randomness they consume grow
faster, so masking overhead is often several-fold and sometimes an order of magnitude.
For lattice schemes the specific difficulty is that they mix \emph{Boolean} operations
(hashing, comparison in the FO check) with \emph{arithmetic} operations (the NTT), and
securely converting between Boolean and arithmetic masking is the expensive, delicate
part.

\paragraph{Blinding.}
A lighter-weight relative is \textbf{blinding}: randomize the \emph{representation} of
a computation without changing its result, so that traces do not line up across runs.
One can add a random multiple of the modulus to a value before an operation, randomize
the projective coordinates in an elliptic-curve step, or shuffle the order in which
independent coefficients are processed so the attacker cannot know which sample in a
trace corresponds to which secret element. Blinding is cheaper than full masking but
weaker; the two are often combined.

\paragraph{Cost and benefit.}
None of this is free, and none of it is always necessary. A server behind a network,
where an attacker can only measure timing remotely, needs constant-time code but rarely
full masking. A smart card or an embedded key store, where the attacker holds the
device and can wire up an oscilloscope, needs masking and fault countermeasures as
well. The right level is a threat-model decision, not a fixed recipe -- but
constant-time coding is the floor beneath all of it.

\section{Pitfalls}

\paragraph{Pitfall 1: assuming the compiler preserves your constant-time code.}
An optimizing compiler is free to rewrite $(a \mathbin{\&} mask)\mathbin{|}(b
\mathbin{\&}\overline{mask})$ back into a branch or a conditional move, or to
short-circuit a masked comparison, if it decides that is faster. Constant-time
behavior is a property of the emitted machine code, not of the source. It must be
checked -- by reading the assembly, using verification tools, or barriers that stop the
optimizer -- not assumed.

\paragraph{Pitfall 2: confusing constant-time with constant wall-clock time.}
Constant-time does not mean the code always runs for the same number of nanoseconds.
It means no \emph{secret}-dependent variation. Variation that depends only on public
data -- for example the length of a message -- is perfectly
acceptable. Chasing literally constant duration is both unnecessary and, on modern
out-of-order CPUs, impossible.

\paragraph{Pitfall 3: reporting a decapsulation failure.}
Returning an error, or logging, or taking a fast path when lattice decryption fails,
hands the attacker exactly the plaintext-checking oracle the FO transform was built to
deny. Decapsulation must always run to completion and always return a key; use implicit
rejection, never an error branch.

\paragraph{Pitfall 4: comparing secrets with \texttt{memcmp}.}
The standard library comparison returns as soon as it finds a differing byte, so its
timing reveals how many leading bytes matched -- fatal for a MAC tag or the FO
ciphertext check. Always use a no-early-exit comparison such as \textsc{ct\_eq}.

\paragraph{Pitfall 5: masking that leaks anyway.}
Two shares that are correct on paper still leak if the hardware or the compiler
recombines them -- for example, if consecutive instructions write both shares to the
same register or bus, or if masking is applied without fresh, independent randomness
each run. Masking is only as good as the independence of its shares in the actual
silicon.

\section{Summary}

It is worth stating the lesson of this chapter plainly. The
post-quantum mathematics we spent the previous chapters building is genuine: it
defeats Shor's algorithm, and no efficient quantum attack on the underlying lattice
and code problems is known. But that mathematics protects only the abstract scheme. A
careless implementation can hand the very same secret to an attacker with an
oscilloscope, a shared cache, or a laser -- an attack that costs minutes and dollars
rather than a quantum computer.

The defense is a discipline, not a gadget:

\begin{itemize}
  \item write \textbf{constant-time} code -- no secret-dependent branches, no
        secret-dependent memory addresses, no variable-time instructions on secrets;
  \item where the threat model includes physical access, add \textbf{masking} and
        \textbf{fault} countermeasures on top;
  \item verify the property in the emitted machine code, because the compiler does not
        promise to preserve it.
\end{itemize}

And this is where the theory of the previous parts closes the loop with practice. The
standards did not treat implementation security as an afterthought; they engineered for
it. ML-KEM's CBD sampler avoids a general discrete Gaussian
(\ref{ch:cbd}); ML-DSA and FN-DSA require careful treatment of secret-dependent
sampling and rejection conditions (\ref{ch:mldsa}, \ref{ch:fndsa}); ML-KEM decapsulation is built on implicit
rejection so success and failure are indistinguishable (\ref{ch:fips203}); and
Falcon's floating-point sampler stands as the standing reminder that when a scheme is
hard to implement safely, that difficulty is itself a security cost (\ref{ch:fndsa}).
Post-quantum security is not only a theorem to be proved. It is a property of the
running system -- and keeping it there, all the way down to the machine code, is not
optional.
